\section{Limitations and practical applicability}\label{sec:limitations}

{\color{blue}\paragraph{Scope of the evidence}
The evaluation covers one year, two sectors of fifteen stocks each, and one market. The evaluation year was not used in any earlier analysis and was evaluated once under a registered protocol, but one year of daily returns limits the precision of Sharpe-ratio comparisons: the 90\% intervals of the registered tests have widths of 0.09--0.23, larger than the differences between forecast sources (at most 0.05). The universes were fixed in hindsight and all firms survived to the end of 2025, which introduces survivorship bias. The paired intervals reflect the resampling of trading days for the 25 agent--seed pairs of each source and condition on one forecaster ensemble per source (the mean of three forecaster seeds), so they do not include the variation that retraining the forecasters with other seeds would add. The forecaster and the policies are defined on a fixed set of stocks and cannot be applied to unseen stocks without retraining. The results therefore describe the behaviour of the framework in these sectors and this period rather than its expected performance elsewhere.

\paragraph{Data and accounting}
Prices are unadjusted and the ledger does not credit cash dividends, which lowers the measured returns of high-dividend stocks such as banks and affects all strategies within a sector alike. The relation graphs are static and use information up to mid-2023, so the forecasts that enter the RL training states (2022--2023) are retrospective with respect to the graphs, whereas the selection and evaluation years are not. The 2025 valuation ratios were rebuilt from a public source to match the definitions of the original panels and agree with them over 2024 (Section~\ref{sec5}); the original panels compute PS and PCF from the previous fiscal year's statements from the first trading day of each year, which anticipates annual reports by up to four months, and the 2025 rows follow the same convention so that the inputs are defined identically in training and evaluation. Execution at the observed close with proportional costs ignores market impact, board lots, and price limits, which matter most for the concentrated forecast-ranking rules.

\paragraph{Model design and market adaptation}
The relation graphs cover 82--92\% of the stock pairs, so they constrain graph attention only weakly; sparser or time-varying graphs, for example from rolling shareholding or news co-occurrence, may separate the encoders more clearly. Because every shareholder edge is also a Kendall edge, the shareholder graph adds no admissible pair in the edge-set formulation; using the relation weights as attention biases or edge features would let the qualitative relations enter the forecaster directly. The forecaster is held fixed during policy learning. Joint fine-tuning with portfolio rewards, cost-aware objectives, longer training histories, and state designs that present the forecast as a ranking rather than as additional rows of a covariance-based state may help the agents use a weak signal; these options were not evaluated and would require the same controlled protocol. Hyperparameters and budgets of the agents were fixed in advance and not tuned, and the agents were not given turnover penalties; the post-hoc analyses of Section~\ref{subsec:posthoc} vary the PPO learning rate and discount factor and the depth of the graph and Transformer components, but not the input window or the learning rate of the forecaster.

\paragraph{Practical use}
The forecaster is inexpensive to train and to run, and its high and low forecasts are consistently more accurate than previous-day values, which may be useful for execution and risk limits. Its tradable-horizon ranking was informative in healthcare in 2025 but varied across years and sectors. In this study, neither the DRL agents nor concentrated ranking rules turned the ranking into a net gain over equal weighting after costs, so practical use would require portfolio construction that limits turnover.\par}

\section{Conclusion}\label{sec6}

{\color{blue}This study proposed MGAT-FinRL, a framework that combines relation graphs built from shareholding, investor-relations, and trading-activity data with a graph-attention Transformer forecaster and DRL portfolio agents, and evaluated it under a protocol registered before the evaluation data were collected. Three findings emerge. First, the forecaster predicts next-day high and low prices 15--29\% more accurately than previous-day values in an unseen year, and its forecasts of the tradable-horizon return carry statistically significant ranking information in healthcare (IC$_2=0.052$, $p=0.012$), but not in banking. Second, the relation graphs add little to this information beyond what a Transformer trained with the same ranking-aware objective obtains without them. Third, DRL agents that observe the forecasts at every decision do not attain a significantly higher Sharpe ratio than agents without forecast information: the difference is 0.043 in healthcare and $-0.040$ in banking, and the learned allocations are nearly insensitive to the forecasts. A simple ranking rule extracts the forecast information before costs but not after proportional costs.

The evaluation therefore locates the constraint of the forecast-to-policy pipeline at the interface between a weak cross-sectional signal and the policy learner, rather than in the forecaster alone. Future work should test policy designs that are built to exploit such signals, including cost-aware objectives and joint training, on longer evaluation periods and broader universes, using the same pre-registered and paired evaluation.\par}
